\section{Filters}

\subsubsection*{Motivation and Intuition}

Fourier series allow us to view a signal as a collection of frequency components. A filter changes a signal by keeping, reducing, or removing selected frequencies.

\begin{conceptbox}
    A filter modifies the Fourier coefficients of a signal. If
    \[
        f(x)=\sum_{n=-\infty}^{\infty}
        c_n e^{i2\pi nx/L},
    \]
    then a filter produces an output
    \[
        y(x)=\sum_{n=-\infty}^{\infty}
        H_n c_n e^{i2\pi nx/L},
    \]
    where $H_n$ is the gain applied to the $n$th harmonic.
\end{conceptbox}

\subsubsection*{Physical Interpretation}

A filter does not usually change all frequencies equally. Instead, it may pass low frequencies, remove high frequencies, isolate a band of frequencies, or remove unwanted oscillations.

\begin{conceptbox}
    Common filters:
    \begin{enumerate}
        \item Low-pass filter: keeps low frequencies and reduces high frequencies.
        \item High-pass filter: keeps high frequencies and reduces low frequencies.
        \item Band-pass filter: keeps frequencies within a selected range.
        \item Band-stop filter: removes frequencies within a selected range.
    \end{enumerate}
\end{conceptbox}

\subsubsection*{Low-Pass Filtering}

A low-pass filter keeps the slow-varying part of a signal. In Fourier series form, this means retaining small values of $|n|$ and removing large values of $|n|$.

For an ideal low-pass filter with cutoff harmonic $N$,
\[
    H_n=
    \begin{cases}
        1, & |n|\leq N, \\
        0, & |n|>N.
    \end{cases}
\]

Thus,
\[
    y(x)=\sum_{n=-N}^{N}c_n e^{i2\pi nx/L}.
\]

\subsubsection*{Worked Example}

\begin{examplebox}
    Suppose a periodic signal has the exponential Fourier series
    \[
        f(x)=3+2\cos\left(\frac{2\pi x}{L}\right)
        +\cos\left(\frac{6\pi x}{L}\right)
        +0.5\cos\left(\frac{20\pi x}{L}\right).
    \]

    Identify the harmonics:
    \[
        3 \quad \text{is the DC component},
    \]
    \[
        2\cos\left(\frac{2\pi x}{L}\right)
        \quad \text{is the first harmonic},
    \]
    \[
        \cos\left(\frac{6\pi x}{L}\right)
        \quad \text{is the third harmonic},
    \]
    \[
        0.5\cos\left(\frac{20\pi x}{L}\right)
        \quad \text{is the tenth harmonic}.
    \]

    Now pass the signal through an ideal low-pass filter that keeps harmonics up to $n=3$.

    The DC, first harmonic, and third harmonic remain. The tenth harmonic is removed.

    Therefore, the filtered output is
    \[
        y(x)=3+2\cos\left(\frac{2\pi x}{L}\right)
        +\cos\left(\frac{6\pi x}{L}\right).
    \]
\end{examplebox}

\subsubsection*{Engineering Insight}

\begin{noteBox}
    Filtering is used in audio processing, communications, power electronics, instrumentation, and data analysis. For example, a rectifier converts AC to a waveform with a DC component and oscillatory components. A low-pass filter can remove the oscillatory harmonics and keep the DC value.
\end{noteBox}

\subsubsection*{Connection to Average Value}

The DC component is the zero-frequency part of a signal. In the exponential Fourier series, it is
\[
    c_0=\frac{1}{L}\int_0^L f(x)\,dx.
\]

In the trigonometric Fourier series, it corresponds to
\[
    \frac{a_0}{2}.
\]

A perfect low-pass filter that removes all oscillatory components leaves only the average value of the signal.

\[
    y(x)=c_0.
\]

\subsubsection*{Common Mistakes and Notes}

\begin{conceptbox}
    Common mistakes:
    \begin{enumerate}
        \item Thinking a filter changes the time domain directly without considering frequency components.
        \item Forgetting that removing high harmonics smooths the waveform.
        \item Forgetting that the DC component is the average value of the signal.
        \item Assuming ideal filters exist physically. Real filters approximate ideal behavior.
    \end{enumerate}
\end{conceptbox}